\subsection{The same whole-section comparison preserves actual tensor products}
The following construction uses the canonical unit and tensor comparisons
for the original sheaf pullback. In particular, its tensor-product identity
is proved from those comparisons rather than prescribed as extra data.

\begin{definition}[Actual whole-section pullback]
\label{def:intrinsic-resume-whole-section-pullback}
\lean{EnvelopeProperty.ClosedProjectiveGlobalSections.wholePullbackSection}
\leanok
For an actual scheme morphism $f:X\to Y$, a module sheaf $A$ on $Y$, and
$s\in H^0(Y,A)$, let $f^\natural s\in H^0(X,f^*A)$ be the pullback
adjunction-unit image of $s$ on $f^{-1}Y$, identified with the equal whole
open $X$.
\end{definition}
\begin{proof}
The pullback adjunction unit is a module morphism from $A$ to the
pushforward of $f^*A$. Its component on the whole target gives a section
on $f^{-1}Y=X$, so it has exactly the asserted source and target.
\end{proof}

\begin{lemma}[The actual unit-normalized value on each whole section]
\label{lem:intrinsic-resume-whole-section-unit-value}
\lean{EnvelopeProperty.ClosedProjectiveGlobalSections.wholePullbackSection_eq_sectionAtOne}
\leanok
\uses{def:intrinsic-resume-whole-section-pullback}
Let $\alpha_s:\mathcal O_Y\to A$ be the original structure-module
morphism with $\alpha_s(1)=s$. Then $f^\natural s$ is the value at $1$
of the actual unit-normalized pullback of $\alpha_s$.
\end{lemma}
\begin{proof}
Naturality of the pullback adjunction unit commutes its component on
$A$ through $\alpha_s$. Its component on the structure module sends the
unit to the inverse actual pullback-unit comparison applied to
$f^\sharp(1)$. The latter is $1$, since the original ring map preserves
units. Evaluating naturality at $1$ therefore gives the stated actual
unit-normalized section value on $X$.
\end{proof}

\begin{lemma}[The same original whole-section isomorphism has this inverse]
\label{lem:intrinsic-resume-whole-section-iso-inverse}
\lean{EnvelopeProperty.ClosedProjectiveGlobalSections.isoPullbackTopSections_inv_wholePullbackSection}
\leanok
\uses{lem:intrinsic-resume-whole-section-unit-value,def:seven-original-iso-pullback-whole-sections-route,thm:seven-iso-top-section-inverse-coherence}
For an actual scheme isomorphism $e:X\simeq Y$, the inverse of the
existing whole-section comparison on $A$ sends every $s$ to
$e^\natural s$.
\end{lemma}
\begin{proof}
Express $s$ as $\alpha_s(1)$ using its original structure-module morphism.
The existing inverse top-section coherence identifies its image with
the value at $1$ of the same unit-normalized pulled morphism. The
preceding lemma identifies that value with $e^\natural s$. Both
identifications hold for every original whole section, proving the
inverse formula for exactly the existing isomorphism.
\end{proof}

\begin{lemma}[Whole pullback commutes with the same module morphisms]
\label{lem:intrinsic-resume-whole-section-pullback-naturality}
\lean{EnvelopeProperty.ClosedProjectiveGlobalSections.wholePullbackSection_map}
\leanok
\uses{def:intrinsic-resume-whole-section-pullback}
For every actual module morphism $a:A\to B$ on $Y$ and every whole
section $s$, the original pulled morphism sends $f^\natural s$ to
$f^\natural(a(s))$.
\end{lemma}
\begin{proof}
Naturality of the pullback adjunction unit with respect to $a$ says
that applying its component on $A$ and then the pushforward of $f^*a$
is the same as first applying $a$ and then its component on $B$.
Evaluate this equality on the whole target and then at $s$. The
identification $f^{-1}Y=X$ gives exactly the stated whole-section equality.
\end{proof}

\begin{theorem}[The canonical tensor pullback multiplies actual whole sections]
\label{thm:intrinsic-resume-whole-section-pullback-tensor}
\lean{EnvelopeProperty.ClosedProjectiveGlobalSections.canonicalTensorPullbackIso_wholePullbackSection}
\leanok
\uses{def:intrinsic-resume-whole-section-pullback,lem:intrinsic-resume-whole-section-pullback-naturality,lem:cartier-ratio-canonical-pullback,lem:tc-seven-canonical-power-transport}
For every scheme morphism $f:X\to Y$, original module sheaves $A,B$ on
$Y$, and whole sections $s,t$, the actual canonical tensor pullback
isomorphism sends $f^\natural(s\otimes t)$ to
$f^\natural s\otimes f^\natural t$.
\end{theorem}
\begin{proof}
Use the genuine coherent tensor product before its canonical comparison
with the explicit sheafified tensor. The pullback cotensorator sends the
adjunction-unit image of a pure tensor section to the pure tensor of the
two unit images. Naturality of the adjunction unit commutes the coherent
to explicit tensor isomorphism through the pullback. In the defining
composite of the canonical explicit tensor pullback, this pulled
isomorphism cancels its inverse. The remaining cotensorator gives the
two pulled factors, and the target coherent to explicit comparison gives
their actual sheafified tensor section. Finally restriction from
$f^{-1}Y$ to the equal whole open commutes with both sheaf morphisms and
the actual tensor-section product, proving the stated whole-section law.
\end{proof}

